\section{Week 8: Alignment, Instruction Tuning, and Preference Optimization}

\subsection{Learning Objectives}

By the end of this week, you should be able to:
\begin{itemize}
  \item Explain why pretraining alone is insufficient for safe deployment
  \item Distinguish capability training from alignment training
  \item Understand supervised instruction tuning (SFT)
  \item Understand preference-based optimization (RLHF / DPO-style)
  \item Identify systems-level and human-factors risks introduced by alignment
\end{itemize}

\subsection{LLM Components Covered}

\textbf{Post-training}
\begin{itemize}
  \item Instruction tuning
  \item Preference learning
  \item Behavior shaping
\end{itemize}

\textbf{Human interface}
\begin{itemize}
  \item Human feedback loops
  \item Normative objectives
  \item Value encoding
\end{itemize}

\subsection{Conceptual Core: Capability vs Alignment}

Pretraining optimizes for:
\begin{quote}
\emph{Statistical prediction of text}
\end{quote}

Alignment optimizes for:
\begin{quote}
\emph{Human-preferred behavior under constraints}
\end{quote}

These objectives are related but not identical.

Alignment does not add new knowledge; it reshapes how existing knowledge is expressed.

\subsection{Supervised Instruction Tuning (SFT)}

SFT fine-tunes a pretrained model on datasets of:
\[
(\text{instruction}, \text{desired response})
\]

Effects:
\begin{itemize}
  \item improves instruction following,
  \item changes tone and verbosity,
  \item reduces irrelevant completions.
\end{itemize}

Limitations:
\begin{itemize}
  \item brittle outside training distribution,
  \item encodes annotator bias,
  \item scales poorly with task diversity.
\end{itemize}

\subsection{Preference-Based Optimization}

Preference optimization uses comparisons:
\[
(\text{prompt}, y^+, y^-)
\]

where $y^+$ is preferred over $y^-$.

Approaches include:
\begin{itemize}
  \item Reinforcement Learning from Human Feedback (RLHF),
  \item Direct Preference Optimization (DPO).
\end{itemize}

These methods optimize relative preference rather than absolute correctness.

\subsection{What Alignment Changes}

Alignment tends to change:
\begin{itemize}
  \item refusal behavior,
  \item politeness and tone,
  \item verbosity,
  \item safety-related responses.
\end{itemize}

It does not:
\begin{itemize}
  \item make models truthful,
  \item eliminate hallucinations,
  \item guarantee ethical behavior.
\end{itemize}

\subsection{Systems Engineering Implications}

\textbf{Requirements}
\begin{itemize}
  \item Predictable behavior under sensitive prompts
  \item Transparency about limitations
\end{itemize}

\textbf{Architecture}
\begin{itemize}
  \item Separation of base model and alignment layers
  \item Auditable post-training pipeline
\end{itemize}

\textbf{Verification \& Validation}
\begin{itemize}
  \item Red-teaming
  \item Behavioral regression testing
\end{itemize}

\textbf{Operations}
\begin{itemize}
  \item Continuous monitoring for drift
  \item Incident response processes
\end{itemize}

\subsection{Human Factors and Failure Modes}

Alignment introduces new human risks:
\begin{itemize}
  \item overtrust due to polite fluency,
  \item false sense of safety,
  \item normative opacity (whose values?).
\end{itemize}

Human users often conflate:
\begin{quote}
Aligned behavior $\neq$ correct or ethical behavior
\end{quote}

\subsection{Key References}

\begin{itemize}
  \item \citet{ouyang2022training} — InstructGPT / RLHF
  \item \citet{rafailov2023direct} — Direct Preference Optimization
\end{itemize}

\subsection{Recommended University Lectures}

\begin{itemize}
  \item Stanford CS324 — alignment and foundation models  
  \url{https://stanford-cs324.github.io/winter2022/}

  \item Berkeley AI Safety lectures  
  \url{https://aisafety.berkeley.edu}
\end{itemize}

\subsection{Practitioner Lab}

See \texttt{labs/week08/week08\_lab\_alignment\_preference.ipynb}.

\subsection{Week 8 Takeaway}

\begin{quote}
Alignment shapes how models behave, not what they fundamentally know.
\end{quote}


